\section{Primitives for quantum scientific computing}
\label{sec:primitives}

The primitive library is the layer that realizes the input-model and
transform strata of \figref{fig:stack}.  We describe its components in the
order in which a solver consumes them: block-encoding algebra
(\secref{sec:be-algebra}), data loading and sparse access
(\secref{sec:dataloading}), the qubitization/QSVT/QSP engine
(\secref{sec:qsvt}), estimation and readout primitives
(\secref{sec:estimation}), reversible arithmetic and the math-function
compiler (\secref{sec:mathfunc}), and the broader scientific-computing suite
(\secref{sec:sc-suite}).  Table~\ref{tab:primitives} summarizes.

\subsection{Block-encoding algebra}
\label{sec:be-algebra}

A \texttt{BlockEncoding} is a first-class object: an operation together with
its normalization factor $\alpha$ (carried as an RIR attribute) and
capability flags.  The library provides the closure operations that
scientific-computing analyses take for granted: products, scalar rescaling,
linear combinations, general LCU assembly (with square-root-weight
PREPAREs), tensor products, adjoints, Kronecker sums, direct sums,
projectors, truncated shift lattices (the structured encodings produced by
finite-difference and stencil discretizations), Pauli words, and small
explicit Pauli expansions for prototyping.  Because every algebraic node is
a real module graph over real oracles, $\alpha$ bookkeeping, capability
propagation, and openness all flow through composition exactly as
\secref{sec:oracles} prescribes (R2, R6).

\subsection{Data loading and sparse access}
\label{sec:dataloading}

Between a mathematical input model and a circuit lies the data-loading
problem, and the library treats it as a family of interchangeable
realizations behind the oracle slots of \secref{sec:qram}:

\begin{itemize}
\item \textbf{Select-swap QROM} with an explicit partition cost model, so
  that the cost of a given table shape can be read off before choosing a
  realization;
\item \textbf{Alias-sampling PREPAREs} for LCU coefficients;
\item \textbf{XOR databases} in gate, QRAM, and banked (multi-bank,
  wide-word) forms;
\item \textbf{Sparse-access pairs}: a location oracle plus an entry XOR
  oracle, in gate and QRAM forms---the input model consumed by
  Childs--Kothari--Somma-style linear-system kernels;
\item \textbf{State preparations}: gate compilations, QRAM-backed
  rotation-angle memories, uniform and basis states, subspace selection, and
  the \emph{application} of a block-encoding to a state.
\end{itemize}

Interchangeability is testable: the example suite runs the same open problem
with gate and QRAM realizations bound to the same slot and cross-checks
full amplitudes on native backends---the semantics must agree bit for bit.

\subsection{Qubitization, QSVT, and a self-verifying QSP engine}
\label{sec:qsvt}

The transform layer is built on qubitization: the library assembles
qubitization walks from PREPARE--SELECT pairs~\cite{low2019qubitization,gilyen2019qsvt}, composes QSVT phase
sequences, and applies oblivious amplitude amplification where the analysis
calls for it.  Its distinctive component is the \emph{quantum signal
processing phase engine}: given a target polynomial, it synthesizes QSP
phases by complement-polynomial rooting (a Durand--Kerner iteration)
followed by layer stripping, supports degrees up to $40$, and
\emph{self-verifies} by round-tripping the synthesized phases to a residual
around $10^{-9}$ before any circuit is built.  Standard transformations are
provided on top: matrix inversion (the $\mathrm{J}$-polynomial response for
QLSS), eigenstate filtering, Hamiltonian simulation via Jacobi--Anger
cosine/sine branches, and fixed-point search in closed form.  This is
currently the most numerically hardened part of the stack: it is exercised
by dedicated circuit-level tests, including amplitude checks on real
backends (\secref{sec:implementation}).

\subsection{Estimation and readout primitives}
\label{sec:estimation}

The estimation family is assembled from the same modularity discipline:
phase estimation~\cite{kitaev1995pea,nielsen2000quantum} keeps its iterate powers as symbolic \texttt{Repeat} nodes
(R5); amplitude estimation~\cite{brassard2002amplitude} is composed as QPE on the Grover iterate with an
explicit decoder; Hadamard tests produce real and imaginary parts; swap
tests compare registers.  All are unitary modules whose \emph{measurement
plans} are standard host-side readouts (\secref{sec:readout}).  Search
infrastructure (Grover iterations, coherent amplification of zero-signal
subspaces) and quantum walks (coined lattice walks and Szegedy walks with
MNRS-style search, driven by pluggable neighbor-table oracles) round out the
algorithmic base layer.

\subsection{Reversible arithmetic and the math-function compiler}
\label{sec:mathfunc}

Scientific-computing circuits eat and produce \emph{numbers}, so arithmetic
is a load-bearing primitive.  The library provides fixed-point reversible
arithmetic---addition, subtraction, negation, absolute value,
multiplication, division, reciprocal, square root, comparisons, selection,
bitwise operations---as \emph{status-flagged} networks: every path sets
domain-invalid and overflow bits that propagate through composition, so a
guarded reciprocal is a first-class object rather than a comment.  Networks
are built from an XOR/AND/NOT SSA representation with compute/uncompute
scheduling over enforced-zero workspaces, and native backends may compile
whole arithmetic modules to C++ operators charged unit cost
(\secref{sec:architecture}).

On top of arithmetic sits the \emph{math-function compiler}, which turns
ordinary pure Python functions into reversible modules without executing
them:

\begin{lstlisting}[style=pyqpython, caption={Compiling a pure Python numerical function to a reversible module (verbatim from the doctested manual; the gas-dynamics pressure function used by the QFVM pipeline of \secref{sec:casestudies}).  The compiler reads the abstract syntax tree (imports restricted to \texttt{math}/\texttt{cmath}; no execution), folds constants, emits a typed SSA middle IR (MIR~0.1), and lowers to an RIR module with fixed-point ports and propagated status flags.  The calling convention is the XOR-update ABI: inputs preserved, result XOR'd into the \texttt{out} register, domain/overflow state XOR'd into \texttt{status}.}, label={lst:mathfunc}]
from pyqecclang import FixedFormat, compile_function, export_toffoli_u3_cz

def pressure(rho, momentum, energy, gamma=1.4):
    velocity = momentum / rho
    return (gamma - 1) * (energy - 0.5 * momentum * velocity)

compiled = compile_function(
    pressure,
    fmt=FixedFormat(width=12, fraction=6),
    constants={"gamma": 1.4},          # generation-time parameter
)
operation = compiled.operation         # ports: rho, momentum, energy,
program = compiled.program()           #       out, status
originir = export_toffoli_u3_cz(program).text
\end{lstlisting}

The frontend performs constant folding and bounded inlining of non-recursive
helpers; the middle IR is a typed SSA DAG (real/complex/bool kinds, explicit
selection of booleans into numerics) with its own schema; the lowering emits
compute / XOR-update / uncompute schedules honoring the workspace ABI.
Elementary kernels (exponentials, trigonometric and hyperbolic functions,
square roots, two-argument arctangent) are polynomial approximations with
per-intrinsic interval bounds, complex kernels decompose into real ones
(e.g.\ $e^{z} = e^{\operatorname{Re}z}(\cos\operatorname{Im}z +
i\sin\operatorname{Im}z)$), integer powers proceed by square-and-multiply,
and domain/accuracy status flags propagate through every path.  The result
is that \emph{classical numerical expertise written as Python} becomes a
reversible oracle implementation---the bridge that applications such as the
Roe flux (\secref{sec:casestudies}) cross.

\subsection{The broader scientific-computing suite}
\label{sec:sc-suite}

Beyond the linear-algebra core, the library includes input models and
algorithms from across quantum scientific computing, each with the same
declare/carry/bind discipline: Heinrich's quantum summation and numerical
integration; Jordan's gradient estimation from phase oracles;
density-matrix input models via \emph{purification access} (with Gibbs-state
preparation by QSVT Chebyshev/Bessel truncation, partial traces, and trace
distance); quantum PCA by density-matrix exponentiation; a quantum SDP
solver scaffold (penalty Hamiltonians, Gibbs solves, trace estimation, and
the matrix-multiplicative-weights outer loop at the host); low-rank
double-factorization and tensor-hypercontraction block-encodings for
quantum chemistry; and decoded quantum interferometry for GF(2)
optimization.  This breadth exists to exercise the language machinery
across input-model variety---and it is what the ODE/PDE stack of the next
section builds on.

\begin{table}[t]
  \centering
  \caption{Primitive inventory, organized by the strata of
  \figref{fig:stack}.  Every family is assembled from modules and honors the
  oracle lifecycle of \secref{sec:oracles}; ``open-slot ready'' means the
  family can be generated against abstract declarations.}
  \label{tab:primitives}
  \small
  \setlength{\tabcolsep}{4pt}\begin{tabular}{@{}p{3.0cm}p{7.1cm}p{3.1cm}@{}}
    \toprule
    \tabhead{Family} & \tabhead{Contents} & \tabhead{Status} \\
    \midrule
    Block-encoding algebra & products, scaling, linear combinations, LCU, tensor, adjoint, Kronecker sum, direct sum, projectors, stencil shifts, Pauli words & implemented; tested; open-slot ready \\
    Data loading & select-swap QROM (partition cost model), alias sampling, XOR databases (gate/QRAM/banked), sparse access pairs, state preparations & implemented; native cross-checked \\
    QSP/QSVT & phase synthesis (self-verifying, degree $\le 40$), qubitization walks, oblivious amplitude amplification, inversion / eigenstate filters / Hamiltonian simulation / fixed-point search & implemented; numerically hardened \\
    Hamiltonian simulation & first-order Pauli Trotter; truncated-Taylor LCU (default kernel); QSP kernel injectable & Trotter/Taylor implemented; general QSP kernel by injection \\
    Estimation & phase estimation (symbolic powers), amplitude estimation + decoder, Hadamard tests, swap tests & implemented \\
    Arithmetic \& math functions & status-flagged fixed-point arithmetic; Python AST $\rightarrow$ SSA MIR $\rightarrow$ RIR compiler with polynomial kernels & implemented; native cross-checked \\
    Scientific suite & integration (Heinrich), gradients (Jordan), density models via purification, Gibbs, QPCA, SDP scaffold, DF/THC low-rank, DQI & implemented at interface level; maturity varies (\secref{sec:implementation}) \\
    \bottomrule
  \end{tabular}
\end{table}
